
\section{Проблема NLIDB и эволюция подходов к Text-to-SQL}

Проблема создания интерфейсов взаимодействия с базами данных на естественном языке (Natural Language Interfaces to Databases, \texttt{NLIDB}) исследуется на протяжении нескольких десятилетий. Основная цель таких систем — преобразовать запрос пользователя в исполняемый \texttt{SQL}-код без потери семантического смысла. Эволюцию подходов к решению задачи \texttt{Text-to-SQL} можно разделить на несколько ключевых этапов, каждый из которых обладает своими преимуществами и ограничениями.

Исторически первыми системами трансляции текста в \texttt{SQL} были анализаторы, основанные на жестко заданных правилах и шаблонах. Системы извлекали ключевые слова из запроса пользователя и сопоставляли их с именами таблиц и столбцов. У этого подхода была крайне низкая масштабируемость. При малейшем изменении схемы базы данных или использовании синонимомов, система не могла корректно обработать запрос.

Далее для решения проблемы масштабируемости стали применяться методы семантического парсинга, где естественный язык сначала переводится в абстрактное логическое представление, а затем компилируется в \texttt{SQL}. Однако эти методы по-прежнему были ограничены в способности обрабатывать сложные запросы и требовали значительных усилий по ручной настройке для каждой новой предметной области. А так же требовался колоссальный объем ручной работы лингвистов и инженеров по созданию грамматик и словарей для каждой новой базы данных, что делало их непрактичными для широкого применения.

И уже с развитием глубокого обучения задача \texttt{Text-to-SQL} стала рассматриваться как задача машинного перевода, где входная последовательность — это текст, а выходная — \texttt{SQL}-код. Использовались рекуррентные нейронные сети (\texttt{RNN}) и архитектуры с механизмами внимания (\texttt{Attention}). Но была так же проблема: склонность к генерации синтаксически неверного кода и проблема «галлюцинаций» — модель могла выдумать несуществующие столбцы или таблицы, не ориентируясь на реальную схему базы данных.

А современный этап развития \texttt{NLIDB} характеризуется использованием предобученных языковых моделей на базе архитектуры \texttt{Transformer}. Эти модели обладают глубоким контекстным пониманием языка и структуры программного кода. Однако для промышленных систем, где критична точность (например, в авиационной аналитике), применение «чистых» генеративных моделей несет риски логических ошибок в сложных запросах с множественными соединениями (JOIN). Поэтому наиболее передовым подходом на сегодняшний день является использование архитектуры \texttt{RAG} (Retrieval-Augmented Generation) в связке с моделями понимания языка.

В рамках данного подхода, схема базы данных, метаданные и примеры сложных запросов векторизуются и помещаются в специализированную векторную базу данных. При поступлении пользовательского запроса система сначала извлекает наиболее релевантный контекст (структуру таблиц, описание специфичных полей). Обогащенный контекст передается в языковую модель для генерации итогового \texttt{SQL}-запроса, валидного для конкретной \texttt{СУБД}.
